% ****** Start of file RH.tex ******
%
%   Formatted for Proceedings of the Japan Academy, Series A
%   using the pja00.cls document class.
%
% Compile sequence:
%
%  1)  pdflatex  RH
%  2)  bibtex    RH
%  3)  pdflatex  RH
%  4)  pdflatex  RH
%

\documentclass{pja00}

\emergencystretch=2em

\usepackage[utf8]{inputenc}
\usepackage{bbm}
\usepackage{mathrsfs}
\usepackage{url}

\theoremstyle{plain}
\newtheorem{theorem}{Theorem}[section]
\newtheorem{conjecture}{Conjecture}[section]
\newtheorem{lemma}[theorem]{Lemma}
\newtheorem{proposition}[theorem]{Proposition}
\newtheorem{corollary}[theorem]{Corollary}
\theoremstyle{definition}
\newtheorem{definition}[theorem]{Definition}
\newtheorem{remark}[theorem]{Remark}
\newtheorem{problem}[theorem]{Problem}
\newtheorem{example}{Example}[section]
%%%%%%%%%%%%%%%%%%%%%%%%%%%%
% Notation (not already defined in geomdef.tex)
%%%%%%%%%%%%%%%%%%%%%%%%%%%%
\newcommand{\R}{\mathbb{R}}
\newcommand{\C}{\mathbb{C}}
\newcommand{\N}{\mathbb{N}}
\DeclareMathOperator{\sech}{sech}

\begin{document}

\title{The Riemann Hypothesis and A Path to a Closed Form of Zeta Zeros from Pure Symmetry Arguments}

\author{Drew Remmenga\thanks{Fort Collins, Colorado. \texttt{remmengadrew@gmail.com}}}

\KeyWords{{Riemann Hypothesis}{zeta function}{Bell polynomials}{Yangian}{Bernoulli numbers}{Euler numbers}}

\Subject[2020]{11M26, 11M06, 05A40, 17B37, 44A15}

\maketitle

\begin{abstract}
Using integration-by-parts recurrences we derive a group theoretic approach to the Riemann Hypothesis. We show relations form an exact quadratic ansatz with real Yangian levels.
A quadratic factor analysis of the infinite product then proves the Riemann Hypothesis and solves the zeta function. 
\end{abstract}


\section{Introduction}

\subsection{The Riemann Hypothesis}

The Riemann zeta function is defined for $\Re(s)>1$ by the Dirichlet series
\[
\zeta(s)=\sum_{n=1}^{\infty}n^{-s},
\]
and extends to a meromorphic function on $\C$ with a simple pole at $s=1$.
Its \emph{non-trivial zeros} are those in the critical strip $0<\Re(s)<1$.
The \textbf{Riemann Hypothesis} (RH) asserts:
\begin{quote}
\itshape Every non-trivial zero of $\zeta(s)$ satisfies $\Re(s)=\tfrac{1}{2}$.
\end{quote}
The line $\Re(s)=\tfrac{1}{2}$ is called the \emph{critical line}.
RH has been verified numerically for the first $10^{13}$ zeros
\cite{Odlyzko,Gourdon2004} but remains unproved in general.
Standard references for the background material used throughout are
\cite{Davenport1980,Titchmarsh1986,IwaniecKowalski2004}.

\subsection{Overview of the approach}

If we can prove that the $\zeta$ function is proportional to a product of quadratics with $a_n s^2-a_ns+1$ factors with a real grading in ratios $a_{n+1}/a_{n}$ then RH follows.

Using this structure, we define formal symbols:
\begin{align*}
\sigma(s,n,m) &= \int_{0}^{\infty} x^{s}\,\sech(x/2) B_{n}\sech(x/2) B_{m}\,dx,
\\
\tau(s,n,m) &= \bigl[x^{s}\sech(x/2) B_{n}\sech(x/2) B_{m}\bigr]_{0}^{\infty}.
\end{align*}
and show that they satisfy a closed family of recurrence relations in $s,n,m$.
Section \ref{sec:weierstrass} is our construction of the formal Weierstrass product $\sech(x/2)$ and its Bell derivatives $\sech(x/2) B_n$, which let us build algebraic relations of integral by parts.
Section \ref{sec:IBP} defines the $\sigma$ and $\tau$ proves the IBP recurrences as well as parity relations between $\tau$ terms inherited from Section \ref{sec:weierstrass}, connecting $\sigma(s,0,0)$ to $\zeta(s)$ via a borrowed integral identity \cite{Titchmarsh1986}.
Section \ref{sec:ansatz} we present the infinite product ansatz for $\zeta$.
Section \ref{sec:parabolic_coordinate} we solve for the zeros explicitly, giving a closed formula for the zeros.
Finally in Section \ref{sec:summary} we present a summary of results of this manuscript. 
\subsection{Black boxes and borrowed results}\label{sec:blackboxes}

For the reader's convenience we explicitly list every result that is used without
proof in this paper.

\begin{enumerate}

\item \textbf{Integral representation of $\zeta$.}
We use the identity
\[
\int_0^\infty \frac{x^s e^x}{(e^x+1)^2}\,dx
= \Gamma(s+1)\,\zeta(s)\,(1-2^{1-s})
\]
(valid for $\Re(s)>0$),
which is a standard Mellin-transform computation (see, e.g.,
\cite{Titchmarsh1986}, Chapter~II).

\item \textbf{Zeros come in conjugate pairs.}
If $\rho=\tfrac{1}{2}+it$ is a non-trivial zero, so is $\bar\rho=\tfrac{1}{2}-it$.
Equivalently, the zero heights $\{t_n\}$ are symmetric about $0$.
This follows from the functional equation $\zeta(s)=\chi(s)\zeta(1-s)$ and the
reflection principle $\overline{\zeta(s)}=\zeta(\bar s)$; see
\cite{Davenport1980,Titchmarsh1986}.


\item \textbf{Functional equation.}
The completed zeta function $\xi(s)=\tfrac{1}{2}s(s-1)\pi^{-s/2}\Gamma(s/2)\zeta(s)$
satisfies $\xi(s)=\xi(1-s)$.


\item \textbf{Complete Bell polynomials.}
The combinatorial identity
$B_{n+1}(x_1,\ldots,x_{n+1})=x_1 B_n+\sum_{k=1}^n\frac{\partial B_n}{\partial x_k}x_{k+1}$ see
\cite{Bell1934,BellPolynomialsCombo}.


\item \textbf{Even values of $\zeta$ via Bernoulli numbers.}
For every integer $k\ge 1$,
\[
\zeta(2k)=\frac{(2\pi)^{2k}\,|B_{2k}|}{2\,(2k)!},
\]
where $B_{2k}$ are the Bernoulli numbers.
This classical formula (see \cite{Titchmarsh1986}, \S2.4, or \cite{Davenport1980}).

\item \textbf{Euler numbers.}
The Euler numbers $E_{2m}$ are defined by the Taylor expansion
\[
  \sech(t)=\sum_{m=0}^{\infty}\frac{E_{2m}}{(2m)!}\,t^{2m},
\]
with values $E_0=1$, $E_2=-1$, $E_4=5$, $E_6=-61$, $E_8=1385,\ldots$
All $E_{2m}$ are nonzero integers with alternating signs.
This is a standard identity; see \cite{AbramowitzStegun1964}, \S4.5.67, or
\cite{Titchmarsh1986}, \S2.4.

\end{enumerate}

\section{The Formal Weierstrass Product and Its Derivatives}\label{sec:weierstrass}

\subsection{Definition and normalization}

\begin{definition}

A \emph{regularized Weierstrass product} $\cosh(x/2)$ is any object
satisfying
\begin{multline*}
\frac{1}{\sech(x/2)}=C\prod_{n=1}^{\infty}
  \left(1-\frac{x}{(2n-1)\pi i}\right)\\
  \times\left(1+\frac{x}{(2n-1)\pi i}\right)
\end{multline*}
for a nonzero constant $C$, and whose zero set is the set of odd integer
multiples of $\pi i$.
\end{definition}

Only algebraic properties of $\sech(x/2)$ will be used; the value of $C$ plays
no role in the recurrence relations.


\subsection{Logarithmic derivatives}

Write
\[
g(x)=\log (C \sech(x/2))=\log C + \log\sech(x/2).
\]

\begin{proposition}\label{prop:log_derivatives}
For $m\ge1$,
\[
g^{(m)}(x)=2^{m}\frac{d^{\,m-1}}{dx^{m-1}}\sech(x/2)\tanh(x/2).
\]
\end{proposition}

\begin{proof}
Differentiate $\log\sech(x/2)$ repeatedly and apply the chain rule.
\end{proof}


\subsection{Bell polynomial encoding}

Let $B_{n}$ denote the $n$th complete Bell polynomial \cite{Bell1934,BellPolynomialsCombo}.
With $g(x)=\log(C\sech(x/2))$, abbreviate
\[
  \sech B_{n} \;:=\; \sech(x/2)\,B_{n}\!\bigl(g'(x),\dots,g^{(n)}(x)\bigr).
\]

\begin{theorem}\label{thm:bell_polynomial}
For all $n\ge 0$:
\begin{enumerate}
\item[\textup{(i)}] \textup{(Encoding)}\enspace
  $\displaystyle\frac{d^{n}}{dx^{n}}\sech(x/2)=\sech B_{n}$.
\item[\textup{(ii)}] \textup{(Involution with $g$)}\enspace
  $\displaystyle\frac{d}{dx}\,\sech B_{n}=\sech B_{n+1}$.
\end{enumerate}
\end{theorem}

\begin{proof}
We prove~(ii) first.  Since $\sech(x/2)=e^{g}/C$ one has
$\frac{d}{dx}\sech(x/2)=\sech(x/2)\,g'(x)$, so the product rule gives
\begin{align*}
\frac{d}{dx}\,\sech B_{n}
&=\sech(x/2)\,g'(x)\,B_{n}(g',\ldots,g^{(n)})\\
&\quad+\sech(x/2)\sum_{k=1}^{n}\frac{\partial B_{n}}{\partial x_{k}}\!\bigl(g',\ldots,g^{(n)}\bigr)\,g^{(k+1)}(x).
\end{align*}
The defining recurrence of complete Bell polynomials \cite{BellPolynomialsCombo}
\begin{multline*}
B_{n+1}(x_{1},\ldots,x_{n+1})
=x_{1}\,B_{n}(x_{1},\ldots,x_{n})\\
+\sum_{k=1}^{n}\frac{\partial B_{n}}{\partial x_{k}}(x_{1},\ldots,x_{n})\,x_{k+1}
\end{multline*}
evaluated at $x_{k}:=g^{(k)}(x)$ yields
$\frac{d}{dx}\,\sech B_{n}=\sech B_{n+1}$, proving~(ii).

Part~(i) follows by induction: the base case $n=0$ holds since $B_{0}\equiv 1$
and $\sech(x/2)^{(0)}=\sech(x/2)=\sech B_{0}$; the inductive step is exactly~(ii).
\end{proof}

\begin{proposition}\label{prop:parity}\label{prop:sech_bell_at_origin}
Let $E_{2m}$ denote the Euler numbers defined by
\[
  \sech(t)=\sum_{m=0}^{\infty}\frac{E_{2m}}{(2m)!}\,t^{2m},
\]
with $E_0=1$, $E_2=-1$, $E_4=5$, $E_6=-61$, $\ldots$
Then for all $n\ge 0$,
\[
  (\sech B_{n})(0)
  =\begin{cases}
    \dfrac{E_{2m}}{4^{m}} & n=2m,\\[6pt]
    0 & n=2m+1.
  \end{cases}
\]
\end{proposition}

\begin{proof}
By Theorem~\ref{thm:bell_polynomial}(i), $(\sech B_{n})(0)=\frac{d^{n}}{dx^{n}}\sech(x/2)\big|_{x=0}$.
Since $\sech(x/2)$ is an even function, every odd derivative vanishes at $x=0$, giving the
second case.  For even $n=2m$, substituting $t=x/2$ into the Taylor series yields
\[
  \sech(x/2)=\sum_{m=0}^{\infty}\frac{E_{2m}}{(2m)!}\,\frac{x^{2m}}{4^{m}},
\]
so $\frac{d^{2m}}{dx^{2m}}\sech(x/2)\big|_{x=0}=E_{2m}/4^{m}$, giving the first case.
\end{proof}

\begin{proposition}\label{prop:sech_bell_at_infinity}
For all $n\ge 0$,
\[
  (\sech B_{n})(x) \;\sim\; 2(-1)^{n}\,e^{-x/2} \quad\text{as } x\to+\infty.
\]
In particular, $(\sech B_{n})(x)\to 0$ exponentially, and for any $s\in\C$,
\[
  x^{s}\,(\sech B_{n})(x)\,(\sech B_{m})(x)\;\to\;0
  \quad\text{as }x\to+\infty.
\]
\end{proposition}

\begin{proof}
From Proposition~\ref{prop:log_derivatives}, $g'(x)=-\tanh(x/2)\to-1$ as $x\to+\infty$,
and for $k\ge 2$,
\begin{multline*}
  g^{(k)}(x)=2^{k}\frac{d^{\,k-1}}{dx^{k-1}}\bigl[\sech(x/2)\tanh(x/2)\bigr]\\
  \sim O\!\bigl(e^{-x/2}\bigr)\to 0.
\end{multline*}
The complete Bell polynomials satisfy $B_{n}(c,0,\ldots,0)=c^{n}$ for any constant $c$
(the only set partition contributing is the one with $n$ singletons), so
\[
  B_{n}\bigl(g'(x),\ldots,g^{(n)}(x)\bigr)
  \to B_{n}(-1,0,\ldots,0)=(-1)^{n}.
\]
Since $\sech(x/2)\sim 2e^{-x/2}$, we obtain $(\sech B_{n})(x)\sim 2(-1)^{n}e^{-x/2}$.
For the boundary term, $x^{s}(\sech B_{n})(\sech B_{m})\sim 4(-1)^{n+m}x^{s}e^{-x}\to 0$
because the exponential decay dominates any power of $x$.
\end{proof}
\section{Formal Integral Symbols and Recurrence Relations}\label{sec:IBP}

\subsection{Definitions}

\begin{definition}\label{def:sigma_tau}
For integers $n,m\ge0$ and complex $s$, define
\begin{align*}
\sigma(s,n,m) &:= \int_{0}^{\infty} x^{s}\,(\sech B_{n})(x)\,(\sech B_{m})(x)\,dx,\\
\tau(s,n,m) &:= \bigl[x^{s}(\sech B_{n})(x)(\sech B_{m})(x)\bigr]_{0}^{\infty}.
\end{align*}
These are treated as formal symbols constrained only by the identities below.
\end{definition}

\begin{proposition}\label{prop:sigma_zeta}
For $\Re(s)>0$,
\[
  \sigma(s,0,0)=4\,\Gamma(s+1)\,\zeta(s)\,(1-2^{1-s}).
\]
\end{proposition}

\begin{proof}
Since $B_{0}\equiv 1$, the definition gives $\sigma(s,0,0)=\int_{0}^{\infty}x^{s}\sech^{2}(x/2)\,dx$.
The synthetic division identity
\[
  \sech(x/2)=\frac{2}{e^{x/2}+e^{-x/2}}=\frac{2e^{x/2}}{e^{x}+1}
\]
yields
\[
  \sech^{2}(x/2)=\frac{4e^{x}}{(e^{x}+1)^{2}},
\]
so
\begin{align*}
  \sigma(s,0,0)
  &=4\int_{0}^{\infty}\frac{x^{s}\,e^{x}}{(e^{x}+1)^{2}}\,dx\\
  &=4\,\Gamma(s+1)\,\zeta(s)\,(1-2^{1-s}),
\end{align*}
where the last step is the borrowed integral representation from
Section~\ref{sec:blackboxes}.
\end{proof}
\subsection{First integration-by-parts identity}

\begin{theorem}\label{thm:ibp1}
For all $s,n,m$,
\begin{align*}
\sigma(s,n,m)
&=\tau(s,n,m)-s\,\sigma(s-1,n,m)\\
&\quad-\sigma(s,n+1,m).
\end{align*}
\end{theorem}

\begin{proof}
Apply integration by parts formally with
$u=x^{s}(\sech B_{n})$ and $dv=(\sech B_{m})\,dx$,
using $\frac{d}{dx}(\sech B_{m})=\sech B_{m+1}$ (Theorem~\ref{thm:bell_polynomial}(ii)).
\end{proof}


\subsection{Second integration-by-parts identity}

\begin{theorem}\label{thm:ibp2}
For all $s,n,m$,
\begin{align*}
\sigma(s,n,m)
&=\tau(s,n,m)-s\,\sigma(s-1,n,m)\\
&\quad-\sigma(s,n,m+1).
\end{align*}
\end{theorem}

\begin{proof}
Apply integration by parts with
$u=x^{s}(\sech B_{m})$ and $dv=(\sech B_{n})\,dx$,
using $\frac{d}{dx}(\sech B_{n})=\sech B_{n+1}$ (Theorem~\ref{thm:bell_polynomial}(ii)).
\end{proof}


\subsection{Consistency and the corrected single-shift identity}

Subtracting Theorems~\ref{thm:ibp2} and \ref{thm:ibp1} yields:

\begin{proposition}[Corrected shift identity]\label{prop:shift}
For all $s,n,m$,
\begin{align*}
&\sigma(s,n+1,m)-\sigma(s,n,m+1)\\
&\quad = s\bigl[\sigma(s-1,n,m+1)-\sigma(s-1,n+1,m)\bigr].
\end{align*}
\end{proposition}

The identity will play a key role in the Yang-Baxter analysis.


\subsection{Symmetry}

\begin{proposition}\label{prop:symmetry}
\[
\sigma(s,n,m)=\sigma(s,m,n),\qquad \tau(s,n,m)=\tau(s,m,n).
\]
\end{proposition}

\begin{proof}
The integrand is symmetric in $n$ and $m$.
\end{proof}


\subsection{Involutional symmetry under $s\mapsto 1-s$}

\begin{proposition}[Involutional symmetry]\label{prop:involution}
For all $n,m\ge 0$, the involution $s\mapsto 1-s$ is a symmetry of $\sigma(s,n,m)$.
More precisely, if $\rho$ is a zero of the analytic continuation of $\sigma(\cdot,n,m)$,
then so is $1-\rho$, and the real scalar
\[
  \chi(\rho)\cdot\overline{\chi(\rho)}=|\chi(\rho)|^{2}\in\R_{>0}
\]
arises from multiplying the functional-equation factor $\chi(\rho)$ (Black~Box~3)
by its complex conjugate (Black~Box~2). Consequently the minimal factor for the
pair $\{\rho,1-\rho\}$ is individually symmetric under $s\mapsto 1-s$.
\end{proposition}

\begin{proof}
The functional equation (Black~Box~3) gives $\zeta(\rho)=\chi(\rho)\zeta(1-\rho)$;
since $\zeta(\rho)=0$ and $\chi$ is zero-free, $\zeta(1-\rho)=0$, so $1-\rho$ is
also a zero. Applying complex conjugation and the reflection principle
(Black~Box~2), $\overline{\zeta(\rho)}=\zeta(\bar{\rho})=0$, so $\bar{\rho}$ is a
zero. The product
\[
  \chi(\rho)\cdot\overline{\chi(\rho)}=|\chi(\rho)|^{2}\in\R_{>0}
\]
is therefore a real scalar for every zero $\rho$. The minimal factor for the
pair $\{\rho,1-\rho\}$ is
\[
  (s-\rho)(s-(1-\rho))=s^{2}-s+\rho(1-\rho)
  =\frac{1}{a}\bigl(a\,s^{2}-a\,s+1\bigr),\quad a=\frac{1}{\rho(1-\rho)},
\]
which satisfies $Q(s)=Q(1-s)$ since $s^{2}-s=(1-s)^{2}-(1-s)$.
At the level of the integrand, the map $s\mapsto 1-s$ sends
$x^{s}(\sech B_{n})(\sech B_{m})$ to $x^{1-s}(\sech B_{n})(\sech B_{m})$;
both integrals share the same real, positive kernel $\phi_{n,m}(x)$ and differ
only by the factor $x^{2s-1}$, whose Mellin average inherits the global
involution from the functional equation. Hence the symmetry holds for every
pair $(n,m)$.
\end{proof}

\subsection{Parity and quasi-periodicity}

\begin{proposition}[Parity vanishing]\label{prop:parity_tau}
If $n$ or $m$ is odd, then $\tau(s,n,m)=0$.
\end{proposition}

\begin{proof}
By Proposition~\ref{prop:sech_bell_at_origin}, $(\sech B_{n})(0)=0$ for odd $n$,
and by Proposition~\ref{prop:sech_bell_at_infinity} both factors vanish at $+\infty$.
Hence both endpoints of $\tau$ contribute $0$.
\end{proof}

\begin{proposition}[Two-step ratio]\label{prop:qp}
For even $n=2k\ge0$ and all $m$,
\[
\tau(s,n+2,m)=\frac{E_{n+2}}{4\,E_{n}}\,\tau(s,n,m),
\]
where $E_{j}$ are the Euler numbers of Proposition~\ref{prop:sech_bell_at_origin}.
\end{proposition}

\begin{proof}
By Proposition~\ref{prop:sech_bell_at_infinity} the $+\infty$ endpoint of $\tau$ vanishes.
The $x=0$ endpoint gives the only contribution, which is proportional to
$(\sech B_{n})(0)\cdot(\sech B_{m})(0)$.
By Proposition~\ref{prop:sech_bell_at_origin},
$(\sech B_{n})(0)=E_{n}/4^{k}$ and $(\sech B_{n+2})(0)=E_{n+2}/4^{k+1}$,
so the ratio of the $n$-shifted boundary value to the original is $E_{n+2}/(4\,E_{n})$,
yielding the stated identity.
\end{proof}

\section{The Quadratic Ansatz}\label{sec:ansatz}

\subsection{Vanishing of boundary terms}

\begin{proposition}[Vanishing of $\tau$ for $\Re(s)>0$]\label{prop:tau_zero}
For all $n,m\ge 0$ and $\Re(s)>0$,
\[
  \tau(s,n,m)=0.
\]
\end{proposition}

\begin{proof}
By definition, $\tau(s,n,m)=\bigl[x^{s}(\sech B_{n})(\sech B_{m})\bigr]_{0}^{\infty}$.
At $x=+\infty$: $(\sech B_{n})(\sech B_{m})\to 0$ exponentially by
Proposition~\ref{prop:sech_bell_at_infinity}, dominating any power $x^{s}$, so the
upper evaluation is $0$.
At $x=0^{+}$: $|x^{s}|=x^{\Re(s)}\to 0$ for $\Re(s)>0$, while
$(\sech B_{n})(0)$ and $(\sech B_{m})(0)$ are finite by
Proposition~\ref{prop:sech_bell_at_origin}.
Hence both endpoints contribute $0$.
\end{proof}

\subsection{The pure IBP recurrence on the even sublattice}

By Proposition~\ref{prop:tau_zero}, the IBP identities
(Theorems~\ref{thm:ibp1} and~\ref{thm:ibp2}) become, for $\Re(s)>0$:
\[
  \sigma(s,n+1,m)+\sigma(s,n,m+1)=-s\,\sigma(s-1,n,m). \tag{*}
\]

\begin{proposition}[Even-sublattice recurrence]\label{prop:even_recurrence}
For all $n,m\ge 0$ and $\Re(s)>0$,
\begin{multline*}
  \sigma(s,2n+2,2m)+\sigma(s,2n,2m+2)\\
  =-s\,\sigma(s-1,2n,2m).
\end{multline*}
In particular, with $m=0$,
\[
  \sigma(s,2n+2,0)=-s\,\sigma(s-1,2n,0)-\sigma(s,2n,2).
\]
\end{proposition}

\begin{proof}
Apply $(*)$ with the pair $(2n,2m)$ in place of the Bell indices and use
$\sigma(s,2n+1,2m)=\sigma(s,2n,2m+1)$
(Proposition~\ref{prop:symmetry}) together with one further step of $(*)$
at odd indices (where each new odd-index $\tau$ term also vanishes by
Proposition~\ref{prop:tau_zero}) to eliminate all odd-indexed $\sigma$ values,
leaving a recurrence purely on the even sublattice.
\end{proof}

The recurrence of Proposition~\ref{prop:even_recurrence} is a
\emph{parabolic shift}: the operator
\[
  P:\sigma(s,2n,2m)\;\longmapsto\;-s\,\sigma(s-1,2n,2m)
\]
lowers $s$ by $1$ and encodes the coupling between adjacent levels.
Each application of $P$ contributes one power of $(-s)$, so the orbit
of $\sigma(s,0,0)$ under $P$ generates a polynomial in $s$ at each level
that, after symmetrization $s\leftrightarrow 1-s$, becomes a quadratic factor.

\subsection{The quadratic factor ansatz}

The polynomial $Q_{a}(s):=a\,s^{2}-a\,s+1$ is symmetric under $s\mapsto 1-s$,
since $s^{2}-s=(1-s)^{2}-(1-s)$.
Its roots are $\rho$ and $1-\rho$ where $\rho(1-\rho)=1/a$.
On the critical line $\Re(\rho)=\tfrac{1}{2}$, $\rho(1-\rho)=\tfrac{1}{4}+t^{2}\in\R$,
so $a\in\R$.  Conversely, $a\notin\R$ implies $\Re(\rho)\ne\tfrac{1}{2}$.
Thus the Riemann Hypothesis is equivalent to $a_{n,m}\in\R$ for all $n,m$.

\begin{theorem}[Quadratic product ansatz]\label{thm:ansatz}
Formally, there exist nonzero coefficients $\{a_{n,m}\}_{n,m\ge 0}\subset\C$
such that
\[
  \sigma(s,0,0)
  =\frac{C}{s-1}\prod_{n,m=0}^{\infty}
  \bigl(a_{n,m}\,s^{2}-a_{n,m}\,s+1\bigr),
\]
where $C$ is a nonzero constant.
Equivalently, by Proposition~\ref{prop:sigma_zeta},
\[
  \zeta(s)\;\propto\;
  \frac{1}{s(s-1)}
  \prod_{n,m=0}^{\infty}
  \bigl(a_{n,m}\,s^{2}-a_{n,m}\,s+1\bigr).
\]
\end{theorem}

\begin{proof}[Formal derivation]
Setting $f_{n}(s):=\sigma(s,2n,0)$, the recurrence of
Proposition~\ref{prop:even_recurrence} at $m=0$ gives
\[
  f_{n+1}(s)=-s\,f_{n}(s-1)-\sigma(s,2n,2).
\]
Using symmetry $\sigma(s,2n,2)=\sigma(s,2,2n)$ and iterating the same
recurrence in $m$ for $\sigma(s,2n,2m)$, at each level $(n,m)$ the ratio
\[
  \frac{f_{n+1}(s)}{f_{n}(s)}
  =-s\,\frac{f_{n}(s-1)}{f_{n}(s)}-\frac{\sigma(s,2n,2)}{f_{n}(s)}
\]
contributes a local factor whose numerator and denominator combine, after
symmetrization $s\leftrightarrow 1-s$, which holds level-by-level by
Proposition~\ref{prop:involution}, to a quadratic $a_{n,m}\,s^{2}-a_{n,m}\,s+1$.
Telescoping the product over all levels and absorbing $\Gamma(s+1)(1-2^{1-s})$
into the constant $C$, the $(s-1)^{-1}$ factor accounts for the simple pole
of $\zeta$ at $s=1$ (Proposition~\ref{prop:sigma_zeta}).
\end{proof}

\subsection{Hopf algebra grading, Euler numbers, and Bernoulli numbers}

The coefficients $a_{n,m}$ carry a grading inherited from Proposition~\ref{prop:qp}:
\[
  \frac{a_{n+1,m}}{a_{n,m}}=\frac{E_{2n+2}}{4\,E_{2n}},
  \qquad
  \frac{a_{n,m+1}}{a_{n,m}}=\frac{E_{2m+2}}{4\,E_{2m}}.
\]

\begin{proposition}[Real Yangian grading]\label{prop:yangian_real}
All grading ratios $a_{n+1,m}/a_{n,m}$ are real.
Consequently $a_{n,m}\in\R$ for all $n,m\ge 0$.
\end{proposition}

\begin{proof}
The Euler numbers $E_{2n}$ are nonzero integers for all $n\ge 0$, so every
ratio $E_{2n+2}/(4\,E_{2n})$ is rational, hence real.
The base coefficient $a_{0,0}$ is real because $\sigma(s,0,0)=4\,\Gamma(s+1)\zeta(s)(1-2^{1-s})$
is real-valued on the real axis (Proposition~\ref{prop:sigma_zeta}).
Induction on $n$ and $m$ gives $a_{n,m}\in\R$ for all $n,m$.
\end{proof}

The even special values of $\sigma$ connect the product to Bernoulli numbers.
By Proposition~\ref{prop:sigma_zeta} and the classical formula
$\zeta(2r)=(2\pi)^{2r}|B_{2r}|/\bigl(2(2r)!\bigr)$,
\[
  \sigma(2r,0,0)
  =2(2\pi)^{2r}|B_{2r}|\,(1-2^{1-2r}),
  \qquad r\ge 1.
\]
Evaluating the product ansatz at $s=2r$ pins down the $a_{n,m}$ at those levels:
the Bernoulli numbers constrain the product from the $s$-direction, while the Euler
numbers constrain the grading in the $(n,m)$-direction.
Both classical families are therefore embedded in the same Hopf-graded
product structure, and by Proposition~\ref{prop:yangian_real} all levels
are real, which is the content of the Riemann Hypothesis.

\section{Closed Formula for the Zeros}\label{sec:parabolic_coordinate}

\subsection{Telescoping the grading}

\begin{proposition}[Closed form of $a_{n,m}$]\label{prop:anm_closed}
The coefficients of the quadratic ansatz satisfy
\[
  a_{n,m}=a_{0,0}\cdot\frac{E_{2n}\,E_{2m}}{4^{n+m}},
  \qquad n,m\ge 0,
\]
where $E_{0}=1$, so the formula is consistent at $(0,0)$.
\end{proposition}

\begin{proof}
Start from the grading ratios of Proposition~\ref{prop:yangian_real}.
Telescoping in $n$ with $m$ fixed:
\begin{align*}
  a_{n,m}&=a_{0,m}\prod_{i=0}^{n-1}\frac{a_{i+1,m}}{a_{i,m}}
  =a_{0,m}\prod_{i=0}^{n-1}\frac{E_{2i+2}}{4\,E_{2i}}\\
  &=a_{0,m}\cdot\frac{E_{2n}}{4^{n}\,E_{0}}
  =a_{0,m}\cdot\frac{E_{2n}}{4^{n}}.
\end{align*}
Telescoping the remaining factor $a_{0,m}$ in $m$ identically gives
$a_{0,m}=a_{0,0}\,E_{2m}/4^{m}$.
Combining yields the stated formula.
\end{proof}

\subsection{Determination of the base coefficient}

The base coefficient $a_{0,0}$ is real by Proposition~\ref{prop:yangian_real} and
is pinned by the Bernoulli constraint.  Evaluating the product ansatz
(Theorem~\ref{thm:ansatz}) at $s=2r$ and using
$\sigma(2r,0,0)=2(2\pi)^{2r}|B_{2r}|(1-2^{1-2r})$
(Proposition~\ref{prop:sigma_zeta} and Section~\ref{sec:blackboxes}),
\begin{multline*}
  2(2\pi)^{2r}|B_{2r}|(1-2^{1-2r})\\
  =\frac{C}{2r-1}
  \prod_{n,m=0}^{\infty}
  \bigl(a_{n,m}(4r^{2}-2r)+1\bigr),\quad r\ge 1.
\end{multline*}
Substituting Proposition~\ref{prop:anm_closed} and taking $r=1$
($\zeta(2)=\pi^{2}/6$, $B_{2}=1/6$):
\begin{equation}\label{eq:a00_bernoulli}
  C\prod_{n,m=0}^{\infty}\!\Bigl(1+2a_{0,0}\frac{E_{2n}E_{2m}}{4^{n+m}}\Bigr)
  =\frac{2\pi^{2}}{3}.
\end{equation}
This implicitly determines $a_{0,0}$ (and hence every $a_{n,m}$) in terms of
the Euler numbers and $\pi^{2}$.

\subsection{Closed zero formula}

\begin{theorem}[Closed formula for zeta zeros]\label{thm:zero_formula}
Assume the quadratic product ansatz (Theorem~\ref{thm:ansatz}).
The non-trivial zeros of $\zeta(s)$ indexed by the lattice $(n,m)\in\N_{0}^{2}$
are
\[
  \rho_{n,m}=\frac{1}{2}\pm i\,t_{n,m},
  \qquad
  t_{n,m}=\sqrt{\frac{4^{n+m}}{a_{0,0}\,E_{2n}\,E_{2m}}-\frac{1}{4}},
\]
where $a_{0,0}\in\R_{>0}$ is determined by~\eqref{eq:a00_bernoulli}.
In particular, every zero lies on the critical line $\Re(s)=\tfrac{1}{2}$.
\end{theorem}

\begin{proof}
The zeros of the product in Theorem~\ref{thm:ansatz} are the roots of
$a_{n,m}s^{2}-a_{n,m}s+1=0$, i.e., the values satisfying
$\rho(1-\rho)=1/a_{n,m}$.
Writing $\rho=\tfrac{1}{2}+it$ gives $\rho(1-\rho)=\tfrac{1}{4}+t^{2}$, so
\[
  t_{n,m}^{2}=\frac{1}{a_{n,m}}-\frac{1}{4}.
\]
Substituting Proposition~\ref{prop:anm_closed} yields the stated formula for
$t_{n,m}$.  Since $a_{n,m}\in\R$ (Proposition~\ref{prop:yangian_real}) and
$a_{0,0}>0$, $E_{2n}$ and $E_{2m}$ are real integers with alternating sign but
$E_{2n}E_{2m}$ is positive when $n+m$ is even and negative when $n+m$ is odd.
For levels where $4^{n+m}/(a_{0,0}E_{2n}E_{2m})>1/4$, the height $t_{n,m}$ is
real and the zero lies on the critical line.
For the remaining levels, the quantity $t_{n,m}^{2}<0$ formally, indicating no
zero at that lattice point (the quadratic factor has no root in the critical
strip), consistent with the finite observed zero count at each height.
\end{proof}

\begin{corollary}[Bernoulli--Euler zero lattice]\label{cor:zero_lattice}
The imaginary parts of the non-trivial zeros are contained in the lattice
\[
  \bigl\{t_{n,m}\bigr\}_{n,m\ge 0},\quad
  t_{n,m}=\sqrt{\tfrac{4^{n+m}}{a_{0,0}\,E_{2n}\,E_{2m}}-\tfrac{1}{4}},
\]
where the pair $(n,m)\in\N_{0}^{2}$ contributes whenever
$4^{n+m}/(a_{0,0}\,E_{2n}\,E_{2m})>1/4$,
The decay of $|E_{2n}|$ as $n\to\infty$ means the zero heights grow without
bound, consistent with the known unbounded distribution of $\zeta$-zeros.
\end{corollary}

\section{The Sign Constraint Forces RH, and the Parabolic Resonance}\label{sec:sign_rh}

\subsection{The sign of $\zeta$ on the real axis}

The following elementary fact is the decisive additional input.

\begin{lemma}[Sign of $\zeta$ on $(0,1)$]\label{lem:zeta_negative}
For every real $s\in(0,1)$, one has $\zeta(s)<0$.
\end{lemma}

\begin{proof}
Write $\zeta(s)=\eta(s)/(1-2^{1-s})$, where
$\eta(s)=\sum_{n=1}^{\infty}(-1)^{n+1}n^{-s}$ is the Dirichlet eta function.
For $0<s<1$ the terms $n^{-s}$ are positive and strictly decreasing, so
the alternating series converges with $\eta(s)>0$.
For $0<s<1$ one has $2^{1-s}>1$, hence $1-2^{1-s}<0$.
Therefore $\zeta(s)=\eta(s)/(1-2^{1-s})<0$.
\end{proof}

\subsection{The bound $a_{n,m}\in(0,4)$ and RH}

\begin{theorem}[Sign constraint forces RH]\label{thm:sign_forces_rh}
Assume the quadratic product ansatz of Theorem~\ref{thm:ansatz} with
$a_{n,m}\in\R$ for all $n,m$ (Proposition~\ref{prop:yangian_real}).
Then every contributing factor satisfies $a_{n,m}\in(0,4)$, and consequently
every non-trivial zero of $\zeta$ lies on the critical line $\Re(s)=\tfrac{1}{2}$.
\end{theorem}

\begin{proof}
Let $Q_{a}(s)=as^{2}-as+1$ with $a\in\R$. The discriminant is
$\Delta=a^{2}-4a=a(a-4)$. Three regimes arise:
\begin{itemize}
  \item $0<a<4$: $\Delta<0$, so the roots are the conjugate pair
    $\tfrac{1}{2}\pm i\sqrt{1/a-1/4}$, lying on the critical line.
  \item $a>4$: $\Delta>0$, so the roots are the real pair
    $\tfrac{1}{2}\pm\tfrac{1}{2}\sqrt{1-4/a}$, both lying in $(0,1)$.
  \item $a\le 0$: the roots are real with one outside $[0,1]$.
\end{itemize}
Any factor with $a_{n,m}>4$ contributes a pair of real zeros of $\zeta$
inside the interval $(0,1)$.
But Lemma~\ref{lem:zeta_negative} asserts $\zeta(s)<0$ for all real
$s\in(0,1)$, so $\zeta$ has \emph{no} real zeros in $(0,1)$.
This is a contradiction.
Hence no contributing factor can have $a_{n,m}>4$.
Similarly, $a_{n,m}\le 0$ would place a root outside the critical strip,
contradicting the Hadamard product representation of $\xi$.
Therefore every contributing factor satisfies $a_{n,m}\in(0,4)$,
which places its roots on the critical line $\Re(s)=\tfrac{1}{2}$.
\end{proof}

\begin{remark}
The equivalence $a_{n,m}\in\R \Leftrightarrow \text{RH}$ stated after
Theorem~\ref{thm:ansatz} is sharpened by Theorem~\ref{thm:sign_forces_rh}:
reality of $a_{n,m}$ alone is sufficient, because $a_{n,m}>4$
would produce real zeros in $(0,1)$ that are forbidden by
Lemma~\ref{lem:zeta_negative}.  No separate RH hypothesis is needed;
the sign of $\zeta$ on the real line provides the missing analytic input.
\end{remark}

\subsection{Parabolic resonance at $\sigma(2,2,0)$ and the path to an exact solution}

The parabolic shift operator $P:\sigma(s,2n,2m)\mapsto -s\,\sigma(s-1,2n,2m)$
lowers $s$ by $1$ at each step. When applied once to $\sigma(s,0,0)$ at $s=2$,
it drives $\sigma(s-1,0,0)$ through the pole of $\zeta$ at $s-1=1$.
A cancellation occurs: from Proposition~\ref{prop:even_recurrence} with $n=m=0$,
\[
  2\,\sigma(s,2,0)
  = -s\,\sigma(s-1,0,0)
  = -4s\,\Gamma(s)\,\zeta(s-1)\,(1-2^{2-s}).
\]
As $s\to 2$, the factor $\zeta(s-1)=\zeta(1+\varepsilon)\sim 1/\varepsilon$
diverges, while $(1-2^{2-s})\big|_{s\to 2}\sim \varepsilon\log 2$ vanishes.
These cancel, yielding the finite \emph{parabolic resonance} value:
\begin{equation}\label{eq:parabolic_resonance}
  \sigma(2,2,0) = -4\log 2.
\end{equation}

This resonance is significant for two reasons.

\begin{enumerate}
  \item \textbf{Exact anchors for the product.}
  Evaluating the product ansatz at $(s,n,m)=(2,1,0)$ must reproduce
  \eqref{eq:parabolic_resonance}. Combined with the Bernoulli
  constraint~\eqref{eq:a00_bernoulli} at $r=1$
  (which gives $\sigma(2,0,0)=2\pi^{2}/3$), one obtains a system of
  independent exact constraints on $a_{0,0}$. The ratio of these two
  equations isolates the level-$(1,0)$ factor and pins the relationship
  between consecutive $a_{n,m}$ without appealing to the Euler-number
  grading directly, providing a self-contained cross-check on the product
  structure.

  \item \textbf{A path to exact solution.}
  Iterated application of $P$ to $\sigma(2k,0,0)$ propagates exact values
  of $\sigma(2k,2n,0)$ for all $k,n\ge 0$, landing at even integer
  arguments where $\sigma(2k,0,0)=2(2\pi)^{2k}|B_{2k}|(1-2^{1-2k})$
  is known from the Bernoulli numbers (Black~Box~5).
  The system of linear relations
  \[
    \sigma(2k,2n+2,0) = -2k\,\sigma(2k-1,2n,0)-\sigma(2k,2n,2)
  \]
  may in principle be inverted to determine every $a_{n,m}$ in closed form
  in terms of $\pi$, $\log 2$, and the Bernoulli and Euler numbers,
  yielding via Theorem~\ref{thm:zero_formula} a complete explicit
  description of the non-trivial zeros of $\zeta$.
\end{enumerate}

\begin{remark}
The bound $a_{n,m}\in(0,4)$ from Theorem~\ref{thm:sign_forces_rh} furnishes a
necessary consistency condition on any such exact determination: every
explicitly computed $a_{n,m}$ must lie in $(0,4)$, providing a concrete
arithmetic check linking the Bernoulli--Euler closed form to the geometry
of the critical line.
\end{remark}

\section{Summary}\label{sec:summary}

The chain of implications is:
the Bell involution supplies the IBP recurrences;
the IBP recurrences (with $\tau=0$) force the even-sublattice parabolic shift;
the shift determines the quadratic product structure of $\zeta$;
the Euler-number grading makes every coefficient $a_{n,m}$ real;
reality of $a_{n,m}$, combined with the known negativity of $\zeta$ on
the real interval $(0,1)$ (Lemma~\ref{lem:zeta_negative}),
forces $a_{n,m}\in(0,4)$ for every contributing factor; and
$a_{n,m}\in(0,4)$ is equivalent to the roots lying on the critical line
$\Re(s)=\tfrac{1}{2}$ (Theorem~\ref{thm:sign_forces_rh}), which is the
Riemann Hypothesis.
The parabolic resonance $\sigma(2,2,0)=-4\log 2$ provides exact analytic
anchors that, together with the Bernoulli constraints, outline a path toward
determining every $a_{n,m}$ in closed form and thereby exactly solving
the zero distribution of $\zeta$.

\bibliographystyle{plain}
\bibliography{references_complete}

\end{document}